\documentclass[10pt,a4paper]{amsart}
\usepackage[margin=19mm]{geometry}
\usepackage{amsmath,amssymb,mathtools}
\usepackage[scr=rsfs,cal=euler]{mathalfa}
\usepackage{xfrac}
\usepackage[unicode,hidelinks,bookmarks=false]{hyperref}
\newcommand{\ie}{i.e.\ }
\newcommand{\cf}{cf.\ }
\newcommand{\resp}{respectively\ }
\newcommand{\Subsection}{\subsection}
\providecommand{\nobreakdash}{\nobreak\hbox{-}}
\setlength{\emergencystretch}{2em}
\multlinegap0pt
\usepackage{amssymb}
\usepackage{enumerate}
\usepackage{tikz}
\usepackage{tcolorbox}
\newtheorem{assumption}{Assumption}
\newtheorem{lemma}{Lemma}
\DeclareMathOperator{\RR}{{\mathbb{R}}}
\DeclareMathOperator{\Xgroundset}{\Upsilon}
\newcommand{\dd}[1]{\,\mathrm{d}{#1}}
\DeclareMathOperator{\imgdom}{{\Omega}}
\DeclareMathOperator{\projvariable}{i}
\DeclareMathOperator{\raygroundset}{R}
\DeclareMathOperator{\rayinverse}{\rho}
\DeclareMathOperator{\rayvariable}{r}
\DeclareMathOperator{\stripfkt}{h}
\title{Determination of Range Conditions for General Projection Pair Operators}
\author{Richard Huber, Rolf Clackdoyle, Laurent Desbat}
\date{Source: arXiv:2510.07480v1 (CC BY 4.0)}
\begin{document}
\maketitle

\noindent\emph{Source and licence.} Determination of Range Conditions for General Projection Pair Operators. Version arXiv:2510.07480v1 (CC BY 4.0), distributed by arXiv under the Creative Commons Attribution 4.0 International licence, http://creativecommons.org/licenses/by/4.0/, as recorded in the arXiv licence metadata for this exact version and shown on the abstract page https://arxiv.org/abs/2510.07480v1. Full licence notice: \texttt{licenses/arxiv-2510.07480-CC-BY-4.0.txt}.\par\smallskip

\noindent\textit{Editorial scope.} Counted unit: the article's lemma lemma\_connected\_separation (source0.tex lines 842--845). The article's complete original proof of it is delivered with it (source0.tex lines 947--961, one \texttt{proof} environment of 15 lines, which the article places away from the statement). No mathematics is written, repaired or re-derived: the statement and the proof are the article's own lines, in the article's own notation, and no retained line is substituted or re-flowed.  Retained uncounted as the source-local context the proof needs: lines 700--703; lines 705--708.  Nothing was omitted from any retained span, so the package carries no omission record.  No retained line contains a citation, so the wrapper carries no bibliography and no bibliography entry.  No retained line contains a figure environment, an image-inclusion command or drawing code, so no image is delivered and the package declares no asset dependency.  Editorial formatting only, in addition: an AMS \texttt{amsart} preamble that restates the article's own declarations and macros used by the retained lines, namely the environments \texttt{assumption}, \texttt{lemma} on separate counters, together with the standard packages those lines need.  The retained environments print in this document as Assumption 1, Lemma 1 in order of appearance, whereas the article prints the same items with the numbering of its own full document.  The complete native source of the article as distributed in the arXiv e-print for this exact version is bundled unchanged as \texttt{sources/186675/source0.tex}.
\begin{assumption}
\label{Assumption_curve_independent}
The map $(\rayinverse_1,\rayinverse_2)\colon \imgdom \to \Xgroundset$ is injective and $\frac{\dd (\rayinverse_1,\rayinverse_2)}{\dd x}$ is regular (i.e., the Jacobi matrix is invertible).
\end{assumption}

\begin{lemma} \label{lemma_X_exists}
Assumption \ref{Assumption_curve_independent} holds if and only if $\Xgroundset$ is open and
there is a continuously differentiable $X\colon \Xgroundset\to \imgdom$ such that $X(\rayinverse_1(x),\rayinverse_2(x))=x$ for $x\in \imgdom$. 
\end{lemma}

\begin{lemma}
\label{lemma_connected_separation}
Let Assumption \ref{Assumption_curve_independent}  hold. When $\stripfkt_{\projvariable}\colon \raygroundset_{\projvariable} \to \RR$ for $\projvariable\in \{1,2\}$ are measurable functions with $\stripfkt_1(\rayinverse_1(x))=\stripfkt_2( \rayinverse_2(x))$ for almost all $x\in \imgdom$, then there is a constant $c\in \RR$ with $\stripfkt_1=c$  and $\stripfkt_2=c$ almost everywhere on $\raygroundset_1$ and $\raygroundset_2$, respectively.
\end{lemma}

\begin{proof}[\textbf{Proof of Lemma \ref{lemma_connected_separation}}]
%We start by showing that $\stripfkt_1$ and $\stripfkt_2$ satisfying $\stripfkt_1(\rayinverse_1(x))=\stripfkt_2( \rayinverse_2(x))$ must be locally constant, from which (via connectedness) we will conclude global constantness.

Let $U_1\subset \raygroundset_1$ and $U_2\subset \raygroundset_2$  be open sets such that $U_1\times U_2\subset \Xgroundset$.
Setting $x=X(\rayvariable_1,\rayvariable_2)$ ($X$ as in Lemma \ref{lemma_X_exists}), we see that 
\begin{equation} \label{equ_proof_assumption_lemma}
 \stripfkt_1(\rayvariable_1)= \stripfkt_1(\rayinverse_1(x))= \stripfkt_2(\rayinverse_2(x))= \stripfkt_2(\rayvariable_2) \qquad \text{ for almost all  $(\rayvariable_1,\rayvariable_2)\in U_1\times U_2$}
 \end{equation}
 (since diffeomorphisms map almost everywhere in $\imgdom$ onto almost everywhere in $\Xgroundset$). Since the far-left  side and the far-right side of \eqref{equ_proof_assumption_lemma} depend only on $\rayvariable_1$ or $\rayvariable_2$, respectively, we see that both sides must be constant almost everywhere. Thus, there is a constant $c\in \RR$ with $\stripfkt_1=c$ on $U_1$ and $\stripfkt_2=c$ and $U_2$.
 
Note that $X$ is surjective, so given any $\rayvariable_1^*\in \raygroundset_1$, there is an $x\in \imgdom$ with $\rayinverse_1(x)=\rayvariable_1^*$ and setting $\rayvariable_2^*:=\rayinverse_2(x)$ we see that $(\rayvariable_1^*,\rayvariable_2^*)\in \Xgroundset$. Since $\Xgroundset$ is open (Lemma \ref{lemma_X_exists}), there are neighborhoods $U_1$ and $U_2$ of $\rayvariable_1^*$ and $\rayvariable_2^*$, respectively, such that $U_1\times U_2\subset \Xgroundset$.
As stated above, $\stripfkt_1$ and $\stripfkt_2$ are constant on such $U_1$ and $U_2$. Thus, $\stripfkt_1$ is a locally constant function, i.e., given any point $\rayvariable_1^* \in \raygroundset_1$, there is an open set $U_1\subset \raygroundset_1$ containing $\rayvariable_1^*$ and a constant $c\in \RR$ such that $\stripfkt_1(\rayvariable_1)=c$ on $U_1$.
A locally constant function $\stripfkt_1$ on an
open and connected set $\raygroundset_1$ is globally constant. The same arguments hold for $\stripfkt_2$.
\end{proof}

\end{document}
